\section{Blind-Issuance Construction}
\label{sec:blind-signature}

This section presents the committed-message issuance protocol underlying ARC-SPRUCE.  The construction combines
the target-shaped commitment of Section~\ref{subsec:prelim-commit}, the BB-tran rational hash of
Section~\ref{subsec:prelim-vsis}, and the SmallWood NIZK-ARK of Section~\ref{subsec:prelim-nizk}.  The semantic
message is \(M=\vecm\Vert\veck\); the DKLW rational-hash characteristic \(\muSig\) is sampled by the issuer.

\subsection{Committed-message target}
The holder commits to the semantic message using
\[
  c=C_MM+A_s s+e.
\]
The issuer samples the BB-tran characteristic and randomness
\[
  \muSig\in\Rq^{\ell_{\mu_{\mathsf{sig}}}},
  \qquad
  x_0\in\Rq^{\ell_{x_0}},
  \qquad
  x_1\in\Rq,
\]
with \(\ell_{\mu_{\mathsf{sig}}}=1\) and \(\ell_{x_0}=2\) for the DKLW BB-tran parameter choice used here.  If
\(B_3-\mathbf 1_nx_1\) is invertible, it sets
\[
  Z=(B_3-\mathbf 1_nx_1)^{-1}
\]
and signs the target
\begin{equation}
  T=B_0+B_1\muSig+B_2x_0+Z+c.
  \label{eq:bs-committed-target}
\end{equation}
The preimage signature is \(\vecu\) with \(\mA\vecu=T\).  This is the IntGenISIS-style shape
\(A\vecu=f(\muSig,\chiSig)+C[M;s;e]\), with \(C=[C_M\mid A_s\mid I]\).

\subsection{Protocol relations}
We record the two relations used later in the issuance and ARC layers.

\paragraph{Pre-sign issuance relation.}
The public statement is
\[
  x_{\mathrm{IssueBS}}=(\ComPublicParam,c),
\]
and the hidden witness is
\[
  w_{\mathrm{IssueBS}}=(M,\vecm,\veck,s,e).
\]
Define \(\mathcal R_{\mathrm{IssueBS}}\) by the constraints
\begin{align}
  c&=C_MM+A_s s+e, \label{eq:issuebs-rel-commit}\\
  M&=\vecm\Vert\veck, \label{eq:issuebs-rel-encode}\\
  \veck&\in\mathcal K_{\mathsf{key}}, \label{eq:issuebs-rel-key}
\end{align}
together with the coefficient bounds on \(M,s,e\) and any issuance policy constraints on \(\vecm\).  The proof only
certifies a well-formed commitment opening and policy/key-space membership.  It does not certify the BB-tran
inverse equation or the final target computation, which are performed by the issuer after the proof is verified.

\paragraph{Post-issuance signature relation.}
For later knowledge statements, it is convenient to isolate the raw committed-message signature relation.  The
public statement is
\[
  x_{\mathrm{SigBS}}=(\mA,B,C_M,A_s,\BoundMsg,\BoundSig),
\]
and the hidden witness is
\[
  w_{\mathrm{SigBS}}=(\vecu,M,s,e,\muSig,x_0,x_1,Z).
\]
Define \(\mathcal R_{\mathrm{SigBS}}\) by
\begin{align}
  (B_3-\mathbf 1_nx_1)\Had Z&=1_n, \label{eq:sigbs-rel-inv}\\
  \mA\vecu&=B_0+B_1\muSig+B_2x_0+Z+C_MM+A_s s+e, \label{eq:sigbs-rel-sig}
\end{align}
with \(M=\vecm\Vert\veck\), \(\veck\in\mathcal K_{\mathsf{key}}\), and the prescribed coefficient bounds on
\(\vecu,M,s,e,\muSig,x_0,x_1\).  No coefficient bound on \(Z\) is imposed unless required by a concrete proof-system
encoding.  The showing relation in Section~\ref{sec:arc-construction} extends this raw signature relation with the
PRF-tag equation.

\subsection{Issuance protocol and raw verification}
\label{subsec:bs-protocol}

\begin{description}
\item[\(\Setup(\SecParam)\):]~
\begin{enumerate}[leftmargin=1.25em]
  \item Generate the commitment matrices \(C_M,A_s\), the message encoding parameters, and distributions
  \(\chi_s,\chi_e\).
  \item Run \(\PublicParamsTrapdoor\gets\TDRH.\Setup(\SecParam,\PublicParamsGen)\) to obtain the BB-tran hash
  description \(B\), the rational-hash sampling spaces, the admissibility error bound \(\delta\), and the
  shortness bound \(\BoundSig\).
  \item Output
  \[
    \PublicParamsBS=(C_M,A_s,\PublicParamsTrapdoor,\BoundMsg,\BoundSig,\mathcal K_{\mathsf{key}}).
  \]
\end{enumerate}

\item[\(\IKeyGen(\PublicParamsBS)\):]~
\begin{enumerate}[leftmargin=1.25em]
  \item Compute \((\mA,\trapdoor)\gets \TDRH.\TrapGen(\PublicParamsTrapdoor)\).
  \item Output \(\vk=\mA\) and \(\sk=\trapdoor\).
\end{enumerate}

\item[\(\Issue\):] The issuer \(\Issuer\) holds \(\sk\) while the holder \(\Holder\) has input semantic message
\(M=\vecm\Vert\veck\) and issuer key \(\vk\).
\begin{enumerate}[leftmargin=1.25em]
  \item The holder samples \(s\leftarrow\chi_s^{k_s}\) and \(e\leftarrow\chi_e^{n_c}\), computes
  \[
    c=C_MM+A_s s+e,
  \]
  and sends \((c,\nizkIssue)\) to the issuer, where
  \[
    \nizkIssue\gets \nizkscheme.\ZKProve
      \bigl(w_{\mathrm{IssueBS}}\mid \mathcal R_{\mathrm{IssueBS}},x_{\mathrm{IssueBS}}\bigr).
  \]
  \item The issuer verifies \(\nizkIssue\).  If it fails, the issuer aborts.
  \item The issuer samples \(\muSig,x_0,x_1\) from the DKLW BB-tran domains.  If
  \(B_3-\mathbf 1_nx_1\notin\Rq^{\times n}\), it rejects this draw and resamples.  The resulting failure
  probability is accounted for as the rational-hash admissibility error.
  \item The issuer computes
  \[
    Z=(B_3-\mathbf 1_nx_1)^{-1},
    \qquad
    T=B_0+B_1\muSig+B_2x_0+Z+c.
  \]
  \item The issuer samples \(\vecu\gets \TDRH.\SampPre(\trapdoor,T)\) and sends
  \((\vecu,\muSig,x_0,x_1)\) to the holder.
  \item The holder checks
  \[
    \mA\vecu=B_0+B_1\muSig+B_2x_0+Z+c
    \qquad\text{and}\qquad
    \|\vecu\|_\infty\le \BoundSig.
  \]
  If the checks succeed, it stores the raw post-issuance witness
  \[
    \sig=(\vecu,M,s,e,\muSig,x_0,x_1).
  \]
\end{enumerate}

\item[\(\Verify(\vk,M,\sig)\):] On input \(\sig=(\vecu,M,s,e,\muSig,x_0,x_1)\), verification rejects if any public
bound or key-space check fails.  It also rejects if \(B_3-\mathbf 1_nx_1\notin\Rq^{\times n}\).  Otherwise it
computes \(Z=(B_3-\mathbf 1_nx_1)^{-1}\) and accepts iff
\[
  \mA\vecu=B_0+B_1\muSig+B_2x_0+Z+C_MM+A_s s+e.
\]
\end{description}

The output \(\sig\) is a raw committed-message witness for the BB-tran signature relation.  The ARC layer uses it
only inside later zero-knowledge showing proofs; the current section does not identify it with a completed blind
signature proof.

\subsection{Correctness and issuer-view hiding}
\begin{proposition}[Correctness]\label{prop:bs-correctness-prop}
Assume completeness of the pre-sign NIZK and correctness of the trapdoor sampler for admissible BB-tran inputs.
Then the issuance layer satisfies Definition~\ref{def:bs-correctness}.
\end{proposition}

\begin{proof}
In an honest execution, the holder samples \((s,e)\), computes \(c=C_MM+A_s s+e\), and produces a complete
pre-sign proof for \(\mathcal R_{\mathrm{IssueBS}}\).  The issuer verifies this proof, samples admissible
\((\muSig,x_0,x_1)\), computes \(Z=(B_3-\mathbf 1_nx_1)^{-1}\), and sets
\[
  T=B_0+B_1\muSig+B_2x_0+Z+c.
\]
Correctness of \(\SampPre\) yields \(\mA\vecu=T\) with \(\|\vecu\|_\infty\le\BoundSig\) except with negligible
failure probability.  The holder and raw verifier recompute the same equation and therefore accept with probability
at least \(1-\negl(\lambda)\).
\end{proof}

\begin{proposition}[Honest-issuer issuer-view hiding]\label{prop:bs-issuer-view-hiding}
Assume Assumption~\ref{ass:mlwe-hiding} for the commitment, zero knowledge of the pre-sign NIZK in the sense of
Definition~\ref{def:nizk-ark}, and honest issuer sampling of \(\muSig,x_0,x_1\) independently of \(M\).  Then the
issuance layer satisfies Definition~\ref{def:bs-blindness} against an honest issuer.
\end{proposition}

\begin{proof}
In the left-right experiment of Definition~\ref{def:bs-blindness}, the issuer first receives \(c\) and
\(\nizkIssue\).  By Lemma~\ref{lem:commit-hiding}, the commitment component for \(M_0\) is computationally
indistinguishable from the commitment component for \(M_1\).  By zero knowledge, the pre-sign proof can be
replaced by a simulated proof depending only on the public statement, incurring only the NIZK simulation advantage.
The values \(\muSig,x_0,x_1,Z\) and the contribution \(B_0+B_1\muSig+B_2x_0+Z\) are generated by the honest issuer
independently of the hidden bit.  Thus the full issuer view in the two worlds is computationally indistinguishable,
and \(\Adv^{\mathsf{ivh}}_{\BlindSignature,\Adv}(\lambda)\) is negligible.
\end{proof}

\begin{remark}[No malicious-issuer blindness theorem]\label{rem:no-malicious-issuer-blindness}
The preceding proposition is an honest-issuer privacy statement.  A malicious-issuer blindness theorem would need a
separate definition and proof, including how adversarial choices of the rational-hash part are handled.
\end{remark}
